%With an initial pressure of $1 bar$ in the evaporation chamber, it is also necessary to validate if the continuum hypothesis remains valid 
%for the air present in the chamber. The Knudsen number relates the mean free path, $\lambda$, to the scale length of the macroscopic flow property gradients 
%and helps determine the validity of the continuum hypothesis. The magnitude of the mean free path can be used to estimate the Knudsen number.
%\begin{equation}
%    \lambda = \frac{\mu}{\rho} \sqrt{\frac{\pi m}{2 k_{B} T}}
%\end{equation}
%For dry air, at a pressure and temperature of $0.03MPa$ and $328$ respectively, molecular mass $m = 4.815e-26\,\text{kg}$, and estimating the kinematic
%viscosity using  Sutherland's formula, the mean free path, $\lambda \approx 2.0732e-06 m$.
\section{Numerical Results}\label{chap3}
This chapter outlines the computational setup and results for simulating the liquid infiltration 
into the evaporation chamber. The employed models are first validated against classical benchmark problems - 
the 1D Stefan problem and the 2D Rayleigh-Taylor instability - to ensure credibility 
of the numerical approach. The Stefan problem 
benchmarks energy conservation and interface motion against the analytical 
solution, and Rayleigh-Taylor simulations reproduce canonical interfacial 
roll-up and mixing trends observed in two-phase flows. 
A concise overview of the methodologies and results for these 
benchmarks are followed by discussion of the infiltration simulations in the experimental geometry.
The infiltration study, built with adaptive remeshing, quantifies free-surface evolution, 
velocity fields, and probe-based pressure/temperature histories in the chamber and nozzle, 
yielding consistent end-of-injection states that can be used to initialise the flashing phase.

\subsection{Stefan Problem}
\label{sec:stef}
Phase change processes are characterised by the motion of an interface 
separating two distinct phases, where heat transfer governs the interface
 displacement according to conservation of energy principles~\cite{Koga2020}.

In the classical Stefan problem, a prescribed temperature boundary propagates a 
melting or freezing front, with the interface movement determined by the balance of latent 
heat and conducted heat~\cite{jonssonOneDimensionalStefan}. Assuming uniform initial 
temperature, constant boundary temperature, and homogeneous thermophysical properties 
in each phase, the semi-infinite one-dimensional Stefan problem represents the simplest 
analytically solvable phase change model. It admits an analytical solution, reducing to 
the standard transient conduction problem when latent heat is neglected, thereby serving 
as a benchmark for verifying phase-change models and the proper implementation of the energy 
equation.~\cite{jonssonOneDimensionalStefan,bolediLevelsetBasedSpacetime2022,andreucciLectureNotesStefan}

In context of this study, the moving-interface heat-balance problem encountered in the 
Stefan problem is analogous to the growth of a spherical vapour bubble in superheated liquid in
late stages. The interface normal velocity is constrained by the interfacial mass flux through density and latent heat -
emphasising the thermal diffusion controlled growth.

\begin{figure}[h]
    \centering
    \includegraphics[width=0.90\textwidth]{stefan_setup}
\caption{Problem setup}
\label{fig:stefan_setup}
\end{figure}
Consider a semi-infinite rectangular domain \(\Omega \subset \mathbb{R}^2\):

\[
\Omega = \{(x,y) \in \mathbb{R} \times \mathbb{R} : x \in [0,\infty), y \in [0, \Delta h]\}\,,
\]

partitioned into solid and liquid subregions \(\Omega_{\ind{s}}\) and \(\Omega_{\ind{l}}\), 
with \(\Omega = \Omega_{\ind{s}} \cup \Omega_{\ind{l}}\). The interface location is 
defined by \(X(t)\), and \(\partial \Omega\) denotes the boundary of \(\Omega\).

The temperature field \(T(x,t)\) in the liquid phase satisfies the heat equation 
coupled with the Stefan condition at the interface. The system reads:
\begin{subequations}\label{eq:stefan}
\begin{align}
\pderiv{T}{t} &= \alpha \ptwoderiv{T}{x}\,, && 0 < x < X(t)\, \quad t_0 < t < t_f, \\
-k \pderiv{T}{x}\Big|_{x = X(t)^-} &= \rho L \deriv{X}{t}\,, && t_0 < t < t_f, \\
T(x,0) &= T_{m}\, && x \in \Omega\, \\
T(0,t) &= T_{l}\, && t_0 < t < t_f\, \\
\nabla T \cdot \vect{n} &= 0\, && \text{on } \partial \Omega \setminus \{x=0\}\,
\end{align}
\end{subequations}

where \(\alpha\) is the thermal diffusivity, \(k\) the thermal 
conductivity, \(\rho\) the density, \(L\) the latent heat, and \(\vect{n}\) 
the outward normal on \(\partial \Omega\).

The one-dimensional analytical solution to \eqref{eq:stefan} is given by:

\begin{equation}
%\begin{align}
T(x,t) = T_l - (T_l - T_m) 
\frac{\operatorname{erf}\left(\frac{x}{2 \sqrt{\alpha t}}\right)}{\operatorname{erf}\left(\frac{X(t)}{2 \sqrt{\alpha t}}\right)}\,, \quad
X(t) = 2 \lambda \sqrt{\alpha t}\,,
%\end{align}
\end{equation}

where \(\lambda\) is determined as the unique root of

\[
f(\lambda) := \frac{C_p (T_m - T_l)}{L} - \frac{\lambda e^{-\lambda^{2}}}{\sqrt{\pi} \operatorname{erf}(\lambda)} = 0\,,
\]

with \(C_p\) the specific heat and \(\operatorname{erf}(\cdot)\) the error function:

\[
\operatorname{erf}(z) := \frac{2}{\sqrt{\pi}} \int_0^z \exp(-y^2) \diff{y}\,.
\]

%--- EDIT: combined and clarified numerical setup with units using the si{} command, reordered for readability ---
The computational domain used to model the Stefan problem in STAR-CCM+
is a two-dimensional rectangular region with dimensions \(L_x = 0.5\,\si{\meter}\) in the 
horizontal direction and \(6 \Delta h = 3.0 \times 10^{-3}\,\si{\meter}\) in the 
vertical direction, where \(\Delta h = 5.0 \times 10^{-4}\,\si{\meter}\) denotes the base mesh size.

A trimmed cell mesher is used to discretize the domain into quadrilateral cells with 
uniform size \(\Delta h\). No mesh refinement or alignment is applied to maintain mesh uniformity.

The simulation models a melting-solidification process using a single-phase, constant-density 
fluid with laminar energy transport. The fluid phase loosely corresponds to water (\(\mathrm{H_2O}\)) with 
constant density \(\rho = 1000\,\si{\kilogram\per\meter^3}\), specific 
heat \(c_p = 4200\,\si{\joule\per\kilogram\per\kelvin}\), and a temperature-dependent 
thermal conductivity switching between solid and liquid values via a user-defined scalar field function. 
The latent heat of fusion is set to \(h_m = 1.0 \times 10^5\,\si{\joule\per\kilogram}\), and 
the melting temperature \(T_m = 273\,\si{\kelvin}\).

The VOF model with a modified High-Resolution Interface Capturing (HRIC) 
convection scheme is employed to track the phase interface. The melting-solidification coupling 
is activated to model latent heat effects during phase change.

Boundary conditions include no-slip velocity constraints on all walls. The upper and 
lower boundaries are adiabatic, while a fixed temperature condition \(T_L = 300\,\si{\kelvin}\) is 
applied on one vertical wall to induce melting. The initial temperature throughout the domain 
is uniform and set to the melting temperature \(T_m = 273\,\si{\kelvin}\), with zero initial velocity.

Solver settings apply the segregated flow solver with SIMPLE pressure-velocity coupling. 
Temporal discretization utilizes an implicit first-order time-stepping scheme with a 
fixed time step \(\Delta t = 0.5\,\si{\second}\). Relaxation factors are set conservatively 
to maintain numerical stability.

Simulation output includes temperature fields, solid volume fraction, and interface location,
 extracted periodically for validation against the analytical solution. The simulation 
 runs until a maximum physical time of \(t_{\mathrm{max}} = 1000\,\si{\second}\).

%This configuration enables accurate modeling of the classical one-dimensional Stefan problem, 
% serving as a validation benchmark for phase change numerical methods and the energy 
% equation discretization in STAR-CCM+.
\begin{table}[h]
\centering
\ra{1.2}
\begingroup
\footnotesize % make only this table use small font
\begin{tabular}{@{}lp{6cm}@{}}
\toprule
\multicolumn{2}{@{}l}{\textbf{Simulation Setup}} \\
\cmidrule(lr){1-2}
Domain & Two-dimensional rectangular \\
Mesh size, \(\Delta h\) & \(2.0 \times 10^{-3}\,\si{\meter}\) \\
Time discretization & Implicit unsteady, first-order Euler \\
Time step, \(\Delta t\) & \(0.5\,\si{\second}\) \\
\addlinespace
\multicolumn{2}{@{}l}{\textbf{Physical Properties}} \\
\cmidrule(lr){1-2}
Density, \(\rho\) & \(1000.0\,\si{\kilogram\per\meter\cubed}\) \\
Dynamic viscosity, \(\mu\) & \(8.8871 \times 10^{-4}\,\si{\pascal\second}\) \\
Specific heat, \(C_p\) & \(4200.0\,\si{\joule\per\kilogram\per\kelvin}\) \\
Latent heat of fusion, \(L\) & \(1.0 \times 10^{5}\,\si{\joule\per\kilogram}\) \\
Thermal conductivity (liquid), \(k_l\) & \(2.33\,\si{\watt\per\meter\per\kelvin}\) \\
Thermal conductivity (solid), \(k_s\) & \(0.62\,\si{\watt\per\meter\per\kelvin}\) \\
Melting temperature, \(T_m\) & \(273\,\si{\kelvin}\) \\
\addlinespace
\multicolumn{2}{@{}l}{\textbf{Numerical Methods}} \\
\cmidrule(lr){1-2}
Multiphase model & VOF with first-order HRIC convection scheme \\
Flow regime & Laminar \\
Discretisation Scheme & Second-order upwind (energy), blended second- and first-order upwind (volume fraction) \\
\bottomrule
\end{tabular}
\endgroup
\caption{Simulation parameters and physics models for the Stefan melting-solidification problem.}
\label{tab:stefan-simulation-params}
\end{table}
Table~\ref{tab:stefan-simulation-params} summarises the key physical properties and numerical methods employed.
The temperature field evolution for a simulation with mesh 
size \(\Delta h = 5.0 \times 10^{-3}\,\si{\meter}\) is shown in \fig{fig:stefan_evolution}. The 
phase interface advances towards the positive \(x\)-direction over time. 
Temperature decreases from the Dirichlet boundary at \(x = 0\) towards the melting 
temperature \(T_m\) at the interface, remaining constant beyond.


\begin{figure}[htbp]
\centering
\includegraphics[width=0.90\textwidth]{stef}
\caption{Evolution of the phase change interface over time.}
\label{fig:stefan_evolution}
\end{figure}

\begin{figure}[htbp]
\centering
\includegraphics[width=0.8\textwidth]{x_PCI}
\caption{Location of the phase change interface versus time.}
\label{fig:stefan_location}
\end{figure}

\begin{figure}[htbp]
\centering
\begin{subfigure}{.49\textwidth}
\includegraphics[width=\linewidth]{err_PCI}
\caption{Absolute error in the location of the interface}
\label{fig:stefan_error_location}
\end{subfigure}
\hfill
\begin{subfigure}{.49\textwidth}
\includegraphics[width=\linewidth]{errL2}
\caption{Temperature profile convergence analysis}
\label{fig:stefan_temp_convergence}
\end{subfigure}
\caption{Error analysis for the Stefan problem across various mesh refinements.}
\label{fig:stefan_convergence}
\end{figure}

\fig{fig:stefan_convergence} compares the numerically predicted interface location for different mesh sizes against the analytical solution. The interface position is inferred by combining volume fraction and temperature data. Data are sampled at \(100\,\si{\second}\) intervals. The absolute interface location error is displayed in \fig{fig:stefan_error_location}. The numerical temperature profiles are quantified against the analytical solution through the \(L_2\) and \(L_\infty\) norms:

\[
L_2 = \sqrt{\frac{1}{N}\sum_{i=1}^{N}(T_i - T_{\mathrm{analytical},i})^2}, \quad
L_\infty = \max_i \left| T_i - T_{\mathrm{analytical},i} \right|.
\]

Temperature error decreases with mesh refinement, as shown in \fig{fig:stefan_temp_convergence}, confirming convergence to the analytical solution.

\subsection{Rayleigh-Taylor Instability}
The \textbf{Rayleigh-Taylor (RT) instability} arises when a denser fluid is accelerated 
into a lighter fluid, causing perturbations at the interface and the formation of 
complex vortical structures~\cite{birkhoffTAYLORINSTABILITYLAMINAR1954}\,. This 
fundamental fluid instability underpins many natural and engineering phenomena. Although 
the idealized, undisturbed RT instability is rarely observed in practice, its core 
mechanisms manifest over a broad range of spatial and temporal scales. These mechanisms 
are influenced by parameters such as viscosity, surface tension, compressibility, 
external accelerations, and stratification~\cite{SHARP19843}\,.

Perturbation growth is driven by opposing gradients of pressure and density localized 
at the interface. The density contrast is quantified by the non-dimensional Atwood number,
\[
At = \frac{\rho_{\ind{H}} - \rho_{\ind{L}}}{\rho_{\ind{H}} + \rho_{\ind{L}}}\,,
\]
where \(\rho_{\ind{H}}\) and \(\rho_{\ind{L}}\) denote the densities of the heavier (upper) and lighter 
(lower) fluids, respectively. Larger Atwood numbers correspond to stronger instability\,.

The nonlinear evolution of the RT instability governs a wide class of two-phase flows 
involving fluid penetration or displacement~\cite{tryggvasonNumericalSimulationsRayleighTaylor1988}\,. 
Simulating this instability accurately is considered a classical benchmark for two-phase 
flow solvers, requiring detailed resolution of interface dynamics and precise implementation 
of interphase momentum and surface tension forces~\cite{gouldingSimplifiedConservativeDiscretization2024}\,.

Consider a two-dimensional rectangular domain \(\Omega \subset \mathbb{R}^2\,\),
\[
\Omega = \{(x,y) \in \mathbb{R} \times \mathbb{R} : x \in [0,L^{*}), y \in [0, 4L^{*}]\}\,
\]
occupied by immiscible, 
viscous, incompressible fluids in subregions \(\Omega_{\ind{H}}\) and \(\Omega_{\ind{L}}\), satisfying 
\(\Omega = \Omega_{\ind{H}} \cup \Omega_{\ind{L}}\). The interface location is described by a scalar height function 
\(h(x,t)\):
\[
\Omega_{\ind{i}}(t) = \left\{ (x,y) \in \mathbb{R}^2 : (-1)^{i}(y - h(x,t)) > 0, \, t \geq 0 \right\}\,, 
\quad \ind{i} \in \{\ind{H}, \ind{L}\}\,,
\]
such that the interface is the curve
\[
X(t) = \left\{ (x,y) \in \mathbb{R}^2 : y = h(x,t),\, t \geq 0 \right\}\,.
\]
\begin{figure}[h ]
    \centering
    \includegraphics[height = 0.80\linewidth]{rt_setup}
\caption{Problem setup}
\label{fig:rt_setup}
\end{figure}
With initial interface \(X(t_0) = X_0\) and initial velocity field 
\(\vect{v}(t_0) = \vect{v}_0\), the governing equations of the Rayleigh-Taylor
instability in the one-fluid formulation are shown in~\eq{eq:rt-equation-system}.
The phases are assumed to be incompressible and the jump conditions at
the interface consist of the kinematic constraint of no relative phase velocity
and the momentum jump.
\begin{subequations}\label{eq:rt-equation-system}
\begin{align}
\pderiv{\rho}{t} + \nabla \cdot (\rho \vect{v}) &= 0 \,,  
&& \text{in } \Omega(t)\,, \\
\pderiv{(\rho \vect{v})}{t} + \nabla \cdot (\rho \vect{v} \otimes \vect{v}) &= -\nabla p 
+ \mu \nabla^{2} \vect{v} + \rho \vect{g} \,,  
&& \text{in } \Omega(t)\,, \\
\nabla \cdot \vect{v} &= 0 \,,  
&& \text{in } \Omega(t)\,, \\
\left[\! \left[\vect{v} \right]\! \right] &= 0 \,,&& \text{on } X(t)\,, \\
\left[\! \left[ \left(
\mu \left(\nabla \vect{v} + (\nabla \vect{v})^{T} \right) - (p + \rho \vect{g} \cdot \vect{x}) \tens{I} 
\right) \cdot \vect{n} \right] \! \right] &= \sigma \kappa \vect{n} + \left[\! \left[\rho \right]\! \right] \vect{g} \,,&& \text{on } X(t)\,, \\
\vect{v}(\vect{x}, t_0) &= \vect{v}_0 \,,  
&& \text{for } t_0 = 0\,, \\
X(t_0) &= X_0 \,,  
&& \text{for } t_0 = 0\,,
\end{align}
\end{subequations}
where \(\vect{n}\) is the unit normal on the interface, \(\tens{I}\) the identity tensor, 
and \(\kappa\) the mean curvature of the interface. 

To impose a single-mode perturbation, the initial interface is defined as
\[
X_0 = 2 L_x + A \cos\left(\frac{2 \pi x}{L_x}\right) \,,
\]
where \(A = 0.1L^{*}\) is the amplitude of the perturbation.

The numerical simulation of the Rayleigh-Taylor instability is performed using STAR-CCM+ 
over a two-dimensional computational domain. The domain, similar to the problem description,  
is a rectangle of size \(L_x = 1.0\,\si{\meter}\) in the horizontal direction 
and \(4 L_x = 4.0\,\si{\meter}\) in the vertical direction.

The domain is discretized using a trimmed cell mesher, producing a uniform 
structured mesh of quadrilateral cells with base size \(\Delta h \si{\meter}\). 
Mesh alignment and local refinements are not applied.

The flow consists of two immiscible, incompressible fluids with constant 
densities \(\rho_{\ind{H}}\) and \(\rho_{\ind{L}}\) and dynamic 
viscosities \(\mu_{\ind{H}}\), \(\mu_{\ind{L}}\). Turbulence is modeled using the Shear Stress 
Transport (SST) \(k\)-\(\omega\) model with low-Reynolds number damping for cases where the 
Reynolds number exceeds the turbulence onset threshold.
The multiphase model is employed with a modified High-Resolution 
Interface Capturing (HRIC) scheme. Surface tension is modeled as a constant 
coefficient \(\sigma = 1.0 \times 10^{-5}\,\si{\newton\per\meter}\), treated 
semi-implicitly. Gravity acts downward as \(\vect{g} = [0,\, -1.0]\,\si{\meter\per\second\squared}\).
The segregated flow solver operates with SIMPLE pressure-velocity coupling. 
Convective and diffusive fluxes are discretized with second-order upwind schemes, and gradients 
are computed via a least squares method. The modified HRIC blends the second- and first-order upwind
schemes while discretising the 
Physical parameters and numerical details are chosen to maintain dynamic 
similarity with the canonical DNS study in~\cite{hamzehlooDirectNumericalSimulations2021}\,. 
Prior literature suggests that a mesh of \(128 \times 512\) cells suffices to capture 
RT instability at \(At \approx 0.75\)~\cite{MIRJALILI2019221}\,. The timestep is 
non-dimensionalized using the Atwood number, and chosen as 
\[
\Delta t = \frac{2.5 \times 10^{-4}}{\sqrt{At}} \, \si{\second}
\]
in line with analogous studies~\cite{DingDiffuseInterfaceModel2007,hamzehlooDirectNumericalSimulations2021}\,.


Boundary conditions consist of slip walls on the top and bottom boundaries and symmetry 
(periodic) boundaries on the left and right sides.
Initial conditions impose a perturbed interface defined by a cosine wave of 
amplitude \(A = 0.1 L_x\,\si{\meter}\), with initial velocity set to zero throughout the domain.

Simulation results are monitored by extracting volume fractions and velocity 
fields along specified planes to analyze interface evolution. 

These simulation parameters ensure dynamic similarity with benchmark direct numerical 
simulations~\cite{hamzehlooDirectNumericalSimulations2021}, providing a reliable 
basis for validation and analysis of Rayleigh-Taylor instability development.
\begin{table}[h]
\centering
\ra{1.2}
\begingroup
\footnotesize % make only this table use small font
\begin{tabular}{@{}lp{6cm}@{}}
\toprule
\multicolumn{2}{@{}l}{\textbf{Simulation Setup}} \\
\cmidrule(lr){1-2}
Mesh type and size, \(\Delta h\) & Quadrilateral trimmed cell mesh, \(4.5 \times 10^{-3}\,\si{\meter}\) \\
Time discretization & Implicit unsteady, first-order Euler \\
Time step, \(\Delta t\) & \(3.5 \times 10^{-3}\,\si{\second}\) \\
\addlinespace
\multicolumn{2}{@{}l}{\textbf{Physical Properties}} \\
\cmidrule(lr){1-2}
Amplitude of perturbation, \(A\) & \(0.1\,\si{\meter}\) \\
Wavelength of perturbation, \(\lambda\) & \(L_x = 1.0\,\si{\meter}\) \\
Surface tension, \(\sigma\) & \(1.0 \times 10^{-5}\,\si{\newton\per\meter}\) \\
Gravity, \(\vect{g}\) & \([0, -1.0]\,\si{\meter\per\second\squared}\) \\
\addlinespace
\multicolumn{2}{@{}l}{\textbf{Numerical Methods}} \\
\cmidrule(lr){1-2}
Multiphase model & VOF with first-order HRIC convection scheme \\
Flow regime & Turbulent, SST \(k\text{-}\omega\) model \\
Discretisation scheme & Second-order upwind (flow, volume fraction) \\
Gradient reconstruction & Least-squares \\
\bottomrule
\end{tabular}
\endgroup
\caption{Simulation parameters and physics models for the Rayleigh-Taylor instability study.}
\label{tab:rt-simulation-params}
\end{table}


\setlength{\textfloatsep}{4pt plus 2.0pt minus 2.0pt}
\begin{figure}[htbp]
    \vspace*{-20pt}
\centering
\begin{subfigure}[b]{\textwidth}
\centering
\includegraphics[scale = 0.67]{rt_re300}
\label{fig:rt_re300}
\end{subfigure}
\vspace*{-0.25pt}
\begin{subfigure}[b]{\textwidth}
\centering
\includegraphics[scale = 0.67]{rt_re1000}
\label{fig:rt_re1000}
\end{subfigure}
\vspace*{-0.25pt}
\begin{subfigure}[b]{\textwidth}
\centering
\includegraphics[scale = 0.67]{rt_re3000}
\label{fig:rt_re3000}
\end{subfigure}
\vspace*{-15pt}
\caption{Evolution of Rayleigh-Taylor instability for Atwood number At = 0.5 at different Reynolds numbers.}
\label{fig:rt_evolution_vertical}
\end{figure}
\setlength{\textfloatsep}{20pt plus 2.0pt minus 2.0pt}
~\fig{fig:rt_evolution_vertical} depicts the evolution of the RT instability for Reynolds numbers 300, 1000, and 3000.
 The initial evolution of the interface is through exponential growth of the infinitesimal disturbance of the interface~\cite{andrewsSmallAtwoodNumber2010}.
 The growth follows linear stability theory which transitions into increasing non-linear behaviour as the amplitude grows to orders of $0.1\lambda$~\cite{SHARP19843}.
 This non-linear evolution is marked by the development of the distinct "bubbles" of lighter fluid rising as the "spikes" of the heavier fluid fall down. 
 The evolving interface experiences increasing vorticity and the increasing velocity shear results in the development of 
 Kelvin-Helmholtz type instabilities~\cite{hamzehlooDirectNumericalSimulations2021,tryggvasonNumericalSimulationsRayleighTaylor1988} leading to the roll-up 
 of the spikes as the amplitude grows to the order of $0.5-1\lambda$.
 For the case of $Re = 3000$, these instabilities are noticeable much earlier all along the interface, leading to a secondary perturbation
 along the boundary of the domain. Further growth leads to secondary vortices along the spike ends which leads to chaotic mixing 
 and complex flow features.~\cite{SHARP19843}.

\subsubsection*{Effect of Reynolds and Atwood Number}\par\vspace{-\baselineskip}\vspace{\parskip}
Increasing the Reynolds number has a relatively small 
 effect on the penetration lengths of the spikes and bubbles - with $Re = 300$ showing the least penetration depths. 
 However, it does have a significant effect on the topology of the interface, with increasingly chaotic vortices visible 
 at the spike rollup and the ensuing turbulent mixing. At $Re = 300$, it is seen in~\fig{fig:rt_evolution_vertical},
 that the flow remains relatively laminar with smooth interfaces, whereas $Re = 3000$ shows significant turbulent mixing
 and vortical complex structures.
\begin{figure}[h!]
\centering
\includegraphics[width=0.98\linewidth]{rt_Reynolds.png}
\caption{Evolution of the RT instability with varying Reynolds and Atwood numbers}
\label{fig:rt_reynolds}
\end{figure}
High Atwood numbers promote non-linearity and increased turbulent mixing, and affect the initial interface growth. 
 The effect and prominence of the counter-rotating vortices at the tip of the spike decreases and reduces 
 the chaotic flow features at the tips of the secondary vortices. The spike also falls much quicker. This is due to the density-driven 
 acceleration outweighing the viscous effects as is seen in~\fig{fig:rt_evolution_high_at} and
 in the very high density ratio simulations of the RT instability in~\cite{HuangConsistentEssentiallyConservative2020}.

\subsubsection*{Grid Convergence}\par\vspace{-\baselineskip}\vspace{\parskip}
In order to evaluate the quality of the mesh, the maximum vertical distance of the 
 spike and bubble tip as measured from the undisturbed initial interface, $y^{*} = |y_{s,b} - y_{0}|$, are evaluated 
 for Reynolds numbers, $Re = 300, 3000$, and Atwood number, $At = 0.5$. Locations of the spike and bubble tips are 
 recorded at every time-step and extrapolated in order to compare to the DNS data available 
 from~\cite{hamzehlooDirectNumericalSimulations2021}.

~\fig{fig:rt_convergence} shows the comparison the results of the simulations with the DNS location of 
 the bubble and spike tips along with the root mean square defined as 
 \[
 err = \displaystyle \sqrt{\sum_{t_0}^{t_{f}} (y_b - y_{dns})^{2}}\,.
 \] 

\begin{figure}[htbp]
\centering
\begin{subfigure}{.49\textwidth}
\includegraphics[width=1.\linewidth]{rt_conv_re300.png}
\caption{$Re = 300$}
\label{fig:rt_conv_300}
\end{subfigure}
\hfill
\begin{subfigure}{.49\textwidth}
\includegraphics[width=1.\linewidth]{rt_conv_re3000.png}
\caption{$Re = 3000$}
\label{fig:rt_conv_3000}
\end{subfigure}
\caption{Grid convergence analysis for Rayleigh-Taylor instability}
\label{fig:rt_convergence}
\end{figure}

\begin{figure}[htbp]
\centering
\begin{subfigure}{.48\textwidth}
\includegraphics[width=0.98\linewidth]{dns_re1000at05.png}
\caption{$Re = 1000$, $At = 0.5$}
\label{fig:dns_re1000}
\end{subfigure}
\hfill
\begin{subfigure}{.48\textwidth}
\includegraphics[width=0.98\linewidth]{dns_re3000at02.png}
\caption{$Re = 3000$, $At = 0.2$}
\label{fig:dns_re3000}
\end{subfigure}
\caption{Comparison with DNS results for different Reynolds numbers and Atwood numbers}
\label{fig:rt_dns_compare}
\end{figure}
\begin{figure}[htbp]
    \centering
    \begin{subfigure}[b]{\textwidth}
        \centering
        \includegraphics[scale = 0.64]{rt_at30}
        \label{fig:rt_at30}
    \end{subfigure}
    \vspace*{-0.25pt}
    \begin{subfigure}[b]{\textwidth}
        \centering
        \includegraphics[scale = 0.64]{rt_at1000}
        \label{fig:rt_at1000}
    \end{subfigure}
    \vspace*{-0.25pt}
    \begin{subfigure}[b]{\textwidth}
        \centering
        \includegraphics[scale = 0.64]{rt_at1500}
        \label{fig:rt_at1500}
    \end{subfigure}
    \caption{Evolution of Rayleigh-Taylor instability for high Atwood numbers At = 30, 1000, and 1500.}
    \label{fig:rt_evolution_high_at}
\end{figure}

~\fig{fig:rt_dns_compare} compares the RT instabilities for $Re = 1000, 3000$ and $At = 0.5, 0.2$ 
 respectively without the evaluation of the errors.

Although, STAR-CCM+, qualitatively, accurately simulates the RT instability, the calculated 
 RMS and absolute errors in the bubble and spike tip locations show erratic behaviour with decreasing mesh size. 
 For $Re = 300$, a sudden increase in the RMS error at $\Delta h = 0.0025$ is seen.
 For higher Reynolds numbers, the bubble tip distances are consistently under-predicted for $Re = 1000, 3000$ at
 later simulation times. For $Re = 3000$, at later simulation times, secondary perturbations develop at the domain edges.
 These secondary perturbations develop into secondary spikes that interfere with the turbulent vortices at the spike tips and the main spike stem. 
 This is, interestingly, not seen for the coarsest mesh size for $Re = 3000$.
 This could be developing Kelvin-Helmholtz instabilities at the surface that transition into secondary RT instabilities~\cite{chenEffectsInitialPerturbations2022}.
 Although Kelvin-Helmholz instabilities develop along the interface at higher Reynolds number, this secondary 
 perturbation is first observed for $Re = 10000$ in~\cite{hamzehlooDirectNumericalSimulations2021}.

This can be ascribed to either the lower order scheme used in the STAR-CCM+ simulations or the choice of 
 turbulence modelling using eddy-viscosity models instead of completely resolving turbulence. For Atwood number 
 that are not close to zero, the Boussinesq hypothesis is also invalid which is an integral assumption in 
 RANS turbulence models. STAR-CCM+, in order to accommodate the spurious currents that may be 
 introduced through the use of eddy-viscosity models, provides models that aid in the damping of momentum and/or 
 turbulence at the interface, however, the addition of these models leads to significant deviations from the DNS results.

\subsection{Liquid Infiltration of Evaporation Chamber}

The evaporation chamber and the nozzle are modeled together as a two-dimensional domain, 
as shown in~\fig{sub1-vorstell}. During the creation of the 2D model, 
STAR-CCM+ prescribes it an automatic and unchangeable depth of \si{1\meter}. To maintain 
kinematic and dynamic similarity with the actual three-dimensional experiment, 
the duration of discharge 
is controlled and the mass flow rates are matched by equating the inlet velocities. 
\begin{figure}
    \centering
    \includegraphics[height = 0.6\linewidth]{Chapters/Imgs/sub2_vof_scene}
    \caption{The meshed domain after the intial adaptive refinement steps for the infiltration problem}
    \label{sub1-vorstell}
\end{figure}
Approximately \si{150\milli\liter} of water is discharged into the chamber during the experiment.
The inlet diameter, \(d_\mathrm{in} = 0.017\si{\meter}\), is taken as the reference length. 
Assuming a discharge time \(t_\mathrm{dis} = 10\si{\second}\) and a constant volumetric flow rate, the inflow velocity of water in the experiment can be calculated as

\[
\dot{Q} = \frac{Q}{t_\mathrm{dis}} = \frac{0.00015}{10} = v \left(\frac{\pi d_\mathrm{in}^2}{4}\right),
\]

where \(Q\) is the total volume discharged. The flow rate is assumed constant over the discharge period. Thus, the inlet velocity depends on the discharge time according to

\[
v = \frac{0.00015}{d_\mathrm{in} \, t_\mathrm{dis}}.
\]
This relation allows the choice of scales to be simple. The diameter of the inlet $d_{\mathrm{in}}$ is taken to be the 
 reference length $L^{*}$ in the simulation. The inflow velocity is thus the reference velocity $U^{*}$, and the resulting reference
 time is given as $t^{*} = L^{*}/U^{*}$.  
The gravitational acceleration vector is assumed to be \(\mathbf{g} = \SI{0.4}{\meter\per\second\squared}\) directed along the negative \(y\)-axis.

In the actual 3D experiment, the water height at the end of the injection under 
steady-state conditions can be estimated, using geometrical data and the volume of 
discharged water, to be approximately \SI{42}{\milli\meter}. In the simulation, 
the water injection is terminated at about \SI{2.75}{\second}, after which the water 
level reaches roughly \SI{42}{\milli\meter}. Subsequently, the simulation continues for 
an additional \SI{5}{\second} with no inflow velocity imposed.

Section probes located along the mid-section of the evaporation chamber (at heights $25, 35,$ and $45\si{\milli\meter}$ respectively) 
and at various 
points along the nozzle measure temperature in a manner replicating the experimental 
measurement setup. The water-air interface is defined by an isosurface corresponding to 
a volume fraction of water equal to 0.5. This isosurface indicates the maximum water 
height along the mid-section probe and is used to record the local 
velocity component in the \(y\)-direction.
Water is modelled as an incompressible liquid whereas the air is modelled as an ideal 
gas. The flow is modelled as unsteady using the volume-of-fluid method and employs a modified HRIC scheme to
capture the interface.

The domain is discretised using a trimmed cell mesher using uniform quadrilateral
 cells refined locally capture interface dynamics accurately.
 The coarsest mesh size is $d_{in}/2$ and the finest mesh size is chosen to be $d_{in}/32$. This refinement
 is performed using the Remeshing solver in conjunction with a user defined field function that
 switches to a finer mesh size based on the local volume fraction. To account for the possibility of
 heat exchange between the system and the surrounding due to conduction and negligible convection,
 the walls of the domain are prescribed using a thermal convection boundary. The coefficient
 of heat transfer is assumed to be $1\si{\watt\per\meter\squared\per\kelvin}$ and the thermal resistance
 is calculated using the chamber dimensions. This resistance is added in parallel to the thermal
 resistance assumed from the insulating foil around the chamber and nozzle. The simulation 
 assumes a total thermal resistance over the system boundary as $0.098\si{\meter\squared\kelvin\per\watt}$.

Surface tension is modeled using a constant coefficient  \(\sigma = 1.0 \times 10^{-5}\,\si{\newton\per\meter}\), treated 
semi-implicitly. Gravity acts downward as \(\vect{g} = [0,\, -0.4]\,\si{\meter\per\second\squared}\).
The segregated flow solver operates with SIMPLE pressure-velocity coupling. 
Convective and diffusive fluxes are discretized with second-order upwind schemes, and gradients 
are computed via a least squares method. The integration in time is performed using the 
first-order implicit Euler scheme. 
\subsubsection*{Spatial and Temporal Convergence}\par\vspace{-\baselineskip}\vspace{\parskip}
    In order to assess the spatial convergence of the solution, four uniform quadrilateral meshes are used. Using the reference length,
    $L^{*}$, the mesh sizes from coarse to fine are $\Delta \text{h}_{1} = d_{in}/4$, $\Delta \text{h}_{2} = d_{in}/8$, $\Delta \text{h}_{3} = d_{in}/16$,
     and $\Delta \text{h}_{4} = d_{in}/32$ respectively. A uniform timestep,
    $\Delta t = 0.001$\,s is used for all the meshes along with the same initial conditions. Data is recorded at every $0.01$\,s.
    The maximum height of the isosurface against the middle section is used to measure the water level. The velocities at 
    various points along the middle section of the chamber are also recorded.~\fig{fig:sub1_conv} depicts the
    recorded water level against simulation time and the y-axis component of the velocities 
    measured along the middle section at the simulation stop time, $t_{f}$ for the above stated mesh sizes. All quantities have been non-dimensionalised
    using the reference length, velocity, and time scales. 
    \begin{figure}[htb]
        \centering
        \includegraphics[width = \textwidth]{sub1_conv}
        \caption{Water level and y-axis velocity component at time, $t_{f}$, for various mesh refinements for timestep $\Delta t = 0.001\,$s.}
        \label{fig:sub1_conv}
    \end{figure}
    \begin{figure}[h]
        \centering
        \includegraphics[width = \textwidth]{sub1_temp_conv}
        \caption{Water level and y-axis velocity component at time, $t_{f}$, for various timestep refinements at $\Delta \text{h}_{4}$.}
        \label{fig:sub1_tconv}
    \end{figure}
    \begin{figure}[htb]
        \centering
        \includegraphics[width = \textwidth]{sub1_mesh_conv}
        \caption{Interface at time, $t = 3.25$\,s, with increasing mesh refinement from left to right - $\Delta \text{h}_{1}\,,\,\Delta \text{h}_{2}\,,\,\Delta \text{h}_{3}\,,$ and $\Delta \text{h}_{4}$.}
        \label{fig:sub1_mesh_conv}
    \end{figure}
    The free surface on the coarsest mesh size, h$_{1}$, tends to behave like a vibrating membrane with no significant development of any waves. The fluid bulk does 
    not contain any air pockets as the mesh is too coarse to resolve any fine details at early times. The finer mesh, h$_{2}$, is able to resolve
    bubbles in the fluid bulk during times of low bubble velocities and when the bubble diameter spans atleast four cells.
    The coarse mesh is, however, able to resolve waves and the oscillatory motion of the isosurface is clearly noticeable.
    The finest mesh size h$_{4}$ is able to resolve merging and breaking apart 
    of the bubbles, droplets escaping the bulk fluid, as well as the flow features developed due to the interaction between the phases at the 
    free surface. The mesh h$_{3}$ while also resolving the bubble interactions in the bulk flow and some coarse droplets, is unable to resolve
    the interface as sharply. However, the increasing void fraction is apparent with decreasing mesh size. 

    The effects of the choice of timestep is similarly evaluated and is seen in~\fig{fig:sub1_tconv}. Using a uniform mesh size $\Delta h_{4}$,
    the simulation is run for timesteps $0.0005\,$s, $0.001\,$s, and $0.002\,$s. With shorter timesteps, the duration 
    of the initial peak in water level gets shorter, however no stark difference in flow behaviour is observed as in the case of spatial 
    refinements. Based on the above observations, the mesh size $\Delta h_{4}$ and timestep $\Delta t = 0.001\,$s is used 
    as the finest mesh and the standard timestep.  
\subsubsection*{Evolution of the Interface and Flow Field}\par\vspace{-\baselineskip}\vspace{\parskip}
The evolution of the solution is highly sensitive even with very conservative values for timesteps and mesh size. The development is 
 mostly symmetric but loses symmetry due to the clashing waves produced by impact against the boundaries. This entails mixing where bubbles
 of air get captured inside the fluid field. Two main churning motions are observed in the fluid flow. A larger circular convection 
 of air is observed in the upper part. This can be due to shear influences of the liquid and the interphase energy exchange. 
\begin{figure}[h]
    \centering 
    \includegraphics[width = \textwidth]{sub1_probe_data}
    \caption{Pressure, temperature, and void fraction measurements using the surface average values at plane sections in the computational domain.}
    \label{fig:sub1_probe_data}
\end{figure}
\begin{figure}[htbp]
    \centering 
    \begin{subfigure}{\linewidth}
    \includegraphics[width = \textwidth]{1}
    \end{subfigure}
    \begin{subfigure}{\linewidth}
    \includegraphics[width = \textwidth]{2}
    \end{subfigure}
    \caption{Flow field visualisations using line integral convolutions and the behaviour of temperature during the infiltration process}
    \label{fig:lefin}
\end{figure}
~\fig{fig:sub1_probe_data} displays the temperature and pressure recorded by the probes in the nozzle area as would be in
the actual experimental setup. The nozzle observes an increasing temperature with increasing water inflow. This temperature 
cools down equally rapidly due to the energy exchange with the water after inflow has stopped. At the end of the simulation, 
the recorded temperatures are lower than the temperatures in the water bulk as shown in the surface averages measured along the 
main body section planes which record a stable temperature of $353\,$K. 
The pressure at the end of the 
infiltration process is recorded to be at $\approx 1.35\si{\bar}$ which is the stipulated pressure at the end of the infiltration
process in the actual experiment. Void fractions measuring existence of air pockets in the water also record almost negligible
values concluding that the infiltration is not turbulent enough to produce thorough mixing of the water and air into a
bubble mixture as the system gets enough time to somewhat stabilise. 
The flow field is visualised using line integral convolutions of the velocity field. Along with this temperature behaviour during infiltration times is depicted in~\fig{fig:lefin}
\subsection{Summary}
Numerical studies of the described benchmarks show convergent Stefan-front 
tracking and qualitatively correct Rayleigh-Taylor instabilit behaviour, 
with noted discrepancies at higher Reynolds numbers attributable to model order and RANS assumptions. 
The infiltration simulations provide resolved interface dynamics and stable 
thermohydraulic end-states (water level, temperature field, and chamber pressure) 
that match the experimental envelope and serve as credible 
initial conditions for subsequent flash-evaporation modeling.